%%
%% This is file `sample-sigconf.tex',
%% generated with the docstrip utility.
%%
%% The original source files were:
%%
%% samples.dtx  (with options: `all,proceedings,bibtex,sigconf')
%% 
%% IMPORTANT NOTICE:
%% 
%% For the copyright see the source file.
%% 
%% Any modified versions of this file must be renamed
%% with new filenames distinct from sample-sigconf.tex.
%% 
%% For distribution of the original source see the terms
%% for copying and modification in the file samples.dtx.
%% 
%% This generated file may be distributed as long as the
%% original source files, as listed above, are part of the
%% same distribution. (The sources need not necessarily be
%% in the same archive or directory.)
%%
%%
%% Commands for TeXCount
%TC:macro \cite [option:text,text]
%TC:macro \citep [option:text,text]
%TC:macro \citet [option:text,text]
%TC:envir table 0 1
%TC:envir table* 0 1
%TC:envir tabular [ignore] word
%TC:envir displaymath 0 word
%TC:envir math 0 word
%TC:envir comment 0 0
%%
%% The first command in your LaTeX source must be the \documentclass
%% command.
%%
%% For submission and review of your manuscript please change the
%% command to \documentclass[manuscript, screen, review]{acmart}.
%%
%% When submitting camera ready or to TAPS, please change the command
%% to \documentclass[sigconf]{acmart} or whichever template is required
%% for your publication.
%%
%%
\documentclass[sigconf]{acmart}

\usepackage{geometry}
\usepackage{graphicx}
\usepackage{hyperref}
\usepackage{natbib}
\usepackage{subcaption}
\usepackage{xcolor}

%%
%% \BibTeX command to typeset BibTeX logo in the docs
\AtBeginDocument{%
  \providecommand\BibTeX{{%
    Bib\TeX}}}

%% Rights management information.  This information is sent to you
%% when you complete the rights form.  These commands have SAMPLE
%% values in them; it is your responsibility as an author to replace
%% the commands and values with those provided to you when you
%% complete the rights form.
%\setcopyright{acmlicensed}
%\copyrightyear{2018}
%\acmYear{2018}
%\acmDOI{XXXXXXX.XXXXXXX}
%% These commands are for a PROCEEDINGS abstract or paper.
%\acmConference[34th ACM SIGSPATIAL]{International Conference on Advances in Geographic Information Systems}{November 3--6, 2026}{Riverside, CA, USA}
\acmConference{}{}{}

%%
%%  Uncomment \acmBooktitle if the title of the proceedings is different
%%  from ``Proceedings of ...''!
%%
%%\acmBooktitle{34th ACM SIGSPATIAL: International Conference on Advances in Geographic Information Systems, November 3--6, 2026, Riverside, CA, USA}
%\acmISBN{978-1-4503-XXXX-X/2018/06}


%%
%% Submission ID.
%% Use this when submitting an article to a sponsored event. You'll
%% receive a unique submission ID from the organizers
%% of the event, and this ID should be used as the parameter to this command.
%%\acmSubmissionID{123-A56-BU3}

%%
%% For managing citations, it is recommended to use bibliography
%% files in BibTeX format.
%%
%% You can then either use BibTeX with the ACM-Reference-Format style,
%% or BibLaTeX with the acmnumeric or acmauthoryear sytles, that include
%% support for advanced citation of software artefact from the
%% biblatex-software package, also separately available on CTAN.
%%
%% Look at the sample-*-biblatex.tex files for templates showcasing
%% the biblatex styles.
%%

%%
%% The majority of ACM publications use numbered citations and
%% references.  The command \citestyle{authoryear} switches to the
%% "author year" style.
%%
%% If you are preparing content for an event
%% sponsored by ACM SIGGRAPH, you must use the "author year" style of
%% citations and references.
%% Uncommenting
%% the next command will enable that style.
%%\citestyle{acmauthoryear}


%%
%% end of the preamble, start of the body of the document source.
\setcopyright{none}
\settopmatter{printacmref=false}
\renewcommand\footnotetextcopyrightpermission[1]{}
% \pagestyle{plain}

\begin{document}

%%
%% The "title" command has an optional parameter,
%% allowing the author to define a "short title" to be used in page headers.
\title[GeoNatureAgent Benchmark]{GeoNatureAgent Benchmark: Benchmarking LLM Agents for Environmental Geospatial Analysis Across Frontier and Open-Weight Foundation Models [Benchmark]}

%%
%% The "author" command and its associated commands are used to define
%% the authors and their affiliations.
%% Of note is the shared affiliation of the first two authors, and the
%% "authornote" and "authornotemark" commands
%% used to denote shared contribution to the research.
\author{Gabriel Diaz-Ireland}
\authornote{Corresponding author.}
\email{gdiazir1@jh.edu}
\orcid{0009-0003-9049-4947}
\affiliation{%
  \institution{Universidad Cat\'{o}lica de \'{A}vila (UCAV)}
  \city{\'{A}vila}
  \country{Spain}
}
\affiliation{%
  \institution{Johns Hopkins University}
  \city{Baltimore, MD}
  \country{USA}
}

\author{Diego Prieto-Herr\'{a}ez}
\email{diego.prieto@ucavila.es}
\orcid{0000-0003-4636-2261}
\affiliation{%
  \institution{Universidad Cat\'{o}lica de \'{A}vila (UCAV)}
  \city{\'{A}vila}
  \country{Spain}}

\author{Mario Garc\'{\i}a Peces}
\email{mario.garcia.peces@gmail.com}
\affiliation{%
  \institution{Independent Researcher}
  \city{Madrid}
  \country{Spain}
}

\author{Javier Vel\'{a}zquez}
\email{javier.velazquez@ucavila.es}
\orcid{0000-0002-9188-3827}
\affiliation{%
  \institution{Universidad Cat\'{o}lica de \'{A}vila (UCAV)}
  \city{\'{A}vila}
  \country{Spain}}

\author{Devika Jain}
\email{kakkar@fas.harvard.edu}
\affiliation{%
  \institution{Center for Geographic Analysis, Harvard University}
  \city{Cambridge, MA}
  \country{USA}
}

%%
%% By default, the full list of authors will be used in the page
%% headers. Often, this list is too long, and will overlap
%% other information printed in the page headers. This command allows
%% the author to define a more concise list
%% of authors' names for this purpose.
\renewcommand{\shortauthors}{Diaz-Ireland et al.}

%%
%% The abstract is a short summary of the work to be presented in the
%% article.
\begin{abstract}
    Environmental scientists spend disproportionate effort on data wrangling rather than analysis. New AI agents can be a helpful tool, but their use still requires validation. Furthermore, no benchmark exists to evaluate AI agents that automate environmental geospatial workflows through structured tool calling against real APIs.
    We introduce the \textbf{GeoNatureAgent Benchmark}, the first benchmark for environmental analysis agents that operate via structured tool calls to a production-style geospatial API. The benchmark comprises 93 tasks across 18 categories, covering municipality-level analysis, multi-turn conversation, spatial reasoning, cross-indicator synthesis, error handling and recovery, ranking, comparison, multilingual understanding, habitat analysis, deep-dive profiling, temporal change, and task rejection.
    Tasks are evaluated against an open, self-hostable geospatial API that serves three environmental indicators across Spain and Portugal via twelve domain-specific tools. We evaluate seven large language models (Claude~Sonnet~4, DeepSeek~V3.2, GLM-5, Gemini~2.5~Pro, Qwen3-235B, GPT-OSS-120B, and Llama~4~Scout) across Vertex~AI and Anthropic platforms under three temperature-1.0 seeds per model, reporting capability and per-case cost as orthogonal axes.
    Results manifest that: (1)~Claude~Sonnet~4 achieves the highest capability at  $60.8\%\pm0.8\%$,  followed closely by DeepSeek~V3.2 at $56.3\%\pm3.1\%$, while no other model exceeds $51\%$; (2)~the cost-accuracy Pareto frontier is occupied entirely by open-weight models (Llama~4~Scout~$\to$~Qwen3-235B~$\to$~DeepSeek~V3.2), with DeepSeek~V3.2 offering $93\%$ of Claude's capability at 11$\times$ lower cost($\$0.011$/case); (3)~comparison tasks remain universally unsolved ($0\%$ across all models for close-value comparisons), exposing systematic reasoning limitations; and (4)~structured tool calling against a real API provides a more discriminative measure of real-world agent capability, with mean accuracies $25$--$35$ percentage points below those reported on general-purpose GIS benchmarks.
    We further demonstrate geographic and domain extensibility by integrating BigEarthNet~V2 land-cover data for Portugal alongside Spanish CO$_2$ suitability and gully erosion indicators.
    GeoNatureAgent Benchmark, the evaluation harness, and the self-hostable API are publicly available\footnote{\url{https://github.com/gabrielireland/GeoNatureAgent_Benchmark}}.
\end{abstract}

%%
%% The code below is generated by the tool at http://dl.acm.org/ccs.cfm.
%% Please copy and paste the code instead of the example below.
%%
\begin{CCSXML}
<ccs2012>
   <concept>
       <concept_id>10010147.10010257.10010293.10010294</concept_id>
       <concept_desc>Computing methodologies~Neural networks</concept_desc>
       <concept_significance>500</concept_significance>
       </concept>
   <concept>
       <concept_id>10002944.10011123.10011130</concept_id>
       <concept_desc>General and reference~Evaluation</concept_desc>
       <concept_significance>500</concept_significance>
       </concept>
   <concept>
       <concept_id>10002951.10003227.10003236.10003237</concept_id>
       <concept_desc>Information systems~Geographic information systems</concept_desc>
       <concept_significance>500</concept_significance>
       </concept>
   <concept>
       <concept_id>10010147.10010178.10010219.10010221</concept_id>
       <concept_desc>Computing methodologies~Intelligent agents</concept_desc>
       <concept_significance>300</concept_significance>
       </concept>
   <concept>
       <concept_id>10010147.10010178.10010219.10010220</concept_id>
       <concept_desc>Computing methodologies~Multi-agent systems</concept_desc>
       <concept_significance>300</concept_significance>
       </concept>
   <concept>
       <concept_id>10010147.10010178.10010179.10003352</concept_id>
       <concept_desc>Computing methodologies~Information extraction</concept_desc>
       <concept_significance>300</concept_significance>
       </concept>
   <concept>
       <concept_id>10010405.10010432.10010437</concept_id>
       <concept_desc>Applied computing~Earth and atmospheric sciences</concept_desc>
       <concept_significance>300</concept_significance>
       </concept>
 </ccs2012>
\end{CCSXML}

\ccsdesc[500]{Computing methodologies~Neural networks}
\ccsdesc[500]{General and reference~Evaluation}
\ccsdesc[500]{Information systems~Geographic information systems}
\ccsdesc[300]{Computing methodologies~Intelligent agents}
\ccsdesc[300]{Computing methodologies~Multi-agent systems}
\ccsdesc[300]{Computing methodologies~Information extraction}
\ccsdesc[300]{Applied computing~Earth and atmospheric sciences}

%%
%% Keywords. The author(s) should pick words that accurately describe
%% the work being presented. Separate the keywords with commas.
\keywords{Geospatial AI, Benchmark, LLM Agents, Tool Calling, Environmental Analysis}

% \received{20 February 2007}
% \received[revised]{12 March 2009}
% \received[accepted]{5 June 2009}

%%
%% This command processes the author and affiliation and title
%% information and builds the first part of the formatted document.
\maketitle

\end{document}
\endinput
%%
%% End of file `sample-sigconf.tex'.
